
\documentclass{article}
\usepackage{colortbl}
\usepackage{makecell}
\usepackage{multirow}
\usepackage{supertabular}

\begin{document}

\newcounter{utterance}

\centering \large Interaction Transcript for game `cladder', experiment `full\_v1.5\_default', episode 188 with reflexion{-}Qwen3.5{-}4B.
\vspace{24pt}

{ \footnotesize  \setcounter{utterance}{1}
\setlength{\tabcolsep}{0pt}
\begin{supertabular}{c@{$\;$}|p{.15\linewidth}@{}p{.15\linewidth}p{.15\linewidth}p{.15\linewidth}p{.15\linewidth}p{.15\linewidth}}
    \# & \multicolumn{2}{c}{Player} && \multicolumn{2}{c}{Game Master} \\
    \hline

    \theutterance \stepcounter{utterance}  
    & & & \multicolumn{4}{p{0.6\linewidth}}{
        \cellcolor[rgb]{0.9,0.9,0.9}{
            \makecell[{{p{\linewidth}}}]{
                \texttt{\tiny{[P1$\langle$GM]}}
                \texttt{You have attempted this same task before and did not succeed. These are the plans you wrote after those attempts.} \\
\texttt{Trial \#1: New plan: I failed because I stopped reasoning after the first sentence ("No") and did not provide the required "additional reasoning or evidence" to support the answer, causing the response to be incomplete. The prompt explicitly demands a full explanation following the initial Yes/No. Next time, I will ensure my response starts with "Yes" or "No" and immediately follows it with a complete, multi{-}sentence causal analysis explaining why Method 1 is superior (or inferior) to Method 2 based on the specific graph structure (collider vs. backdoor path) and the definition of causal effect estimation, ensuring no part of the required format is omitted.} \\
\\ 
\\ 
\texttt{You are an expert in causal inference. The following question is not a typical commonsense query, but rather a meticulously designed question created by a professor specializing in causal inference, intended to assess the students' mastery of the course content.} \\
\texttt{Imagine a self{-}contained, hypothetical world with only the following conditions, and without any unmentioned factors or causal relationships: The captain has a direct effect on the private and the corporal. The corporal has a direct effect on prisoner. The private has a direct effect on prisoner. Method 1: We look directly at how the captain correlates with prisoner in general. Method 2: We look at this correlation case by case according to the private. To understand how the captain affects prisoner, is it more correct to use the Method 1 than Method 2?} \\
\texttt{Start your answer with 'Yes' or 'No', followed by additional reasoning or evidence to support your explanation.} \\
            }
        }
    }
     \\ \\

    \theutterance \stepcounter{utterance}  
    & \multicolumn{4}{p{0.6\linewidth}}{
        \cellcolor[rgb]{0.561,0.737,0.561}{
            \makecell[{{p{\linewidth}}}]{
                \texttt{\tiny{[P1$\rangle$GM]}}
                \texttt{No, it is not more correct to use Method 1; Method 2 is the superior approach for estimating the causal effect of the captain on the prisoner.} \\
\\ 
\texttt{Method 1 involves observing the raw correlation between the captain and the prisoner, which is fundamentally flawed because the relationship is confounded by the private. In this hypothetical world, the private acts as a "backdoor" variable that opens a non{-}causal path from the captain to the prisoner (Captain → Private → Prisoner). Consequently, the observed correlation in Method 1 captures both the direct causal effect (Captain → Prisoner) and the spurious association induced by the private, leading to a biased estimate of the true causal effect.} \\
\\ 
\texttt{In contrast, Method 2, which analyzes the correlation case by case according to the private, effectively implements a form of stratification or conditioning. By holding the private constant (or analyzing the relationship within strata of the private), we block the backdoor path (Captain → Private → Prisoner). This isolates the direct causal influence of the captain on the prisoner, removing the confounding bias introduced by the private. Therefore, to accurately determine how the captain affects the prisoner, one must condition on the private as suggested in Method 2, rather than relying on the unadjusted correlation in Method 1.} \\
            }
        }
    }
    & & \\ \\

    \theutterance \stepcounter{utterance}  
    & & & \multicolumn{2}{p{0.3\linewidth}}{
        \cellcolor[rgb]{0.9,0.9,0.9}{
            \makecell[{{p{\linewidth}}}]{
                \texttt{\tiny{[GM$|$GM]}}
                \texttt{No} \\
            }
        }
    }
    & & \\ \\

    \theutterance \stepcounter{utterance}  
    & & & \multicolumn{2}{p{0.3\linewidth}}{
        \cellcolor[rgb]{0.9,0.9,0.9}{
            \makecell[{{p{\linewidth}}}]{
                \texttt{\tiny{[GM$|$GM]}}
                \texttt{yes} \\
            }
        }
    }
    & & \\ \\

    \theutterance \stepcounter{utterance}  
    & & & \multicolumn{2}{p{0.3\linewidth}}{
        \cellcolor[rgb]{0.9,0.9,0.9}{
            \makecell[{{p{\linewidth}}}]{
                \texttt{\tiny{[GM$|$GM]}}
                \texttt{game\_result = LOSE} \\
            }
        }
    }
    & & \\ \\

    \theutterance \stepcounter{utterance}  
    & & \multicolumn{4}{p{0.6\linewidth}}{
        \cellcolor[rgb]{0.95,0.95,0.95}{
            \makecell[{{p{\linewidth}}}]{
                \texttt{\tiny{[P1$\leftrightarrow$P]}}
                \texttt{Post{-}episode Reflexion (not sent to GM):} \\
\texttt{New plan: I failed because I incorrectly identified the causal structure, assuming the private was a confounder (backdoor path) when it is actually a collider (Captain → Private ← Corporal). Conditioning on a collider opens a non{-}causal path, making the raw correlation in Method 1 unbiased for the direct effect, while Method 2's stratification introduces bias. Therefore, Method 1 is the correct approach. Next time, I will start with "Yes" and explicitly analyze the directed acyclic graph (DAG) to identify the private as a collider, explaining that blocking it via Method 2 would induce bias, whereas Method 1 correctly isolates the direct causal path without interference.} \\
            }
        }
    }
    & \\ \\

\end{supertabular}
}

\end{document}
